\chapter{CLI e Usabilità}

\section{Framework Click}

L'interfaccia a riga di comando è costruita con \textbf{Click}, un framework Python che offre:

\begin{itemize}
    \item Parsing automatico degli argomenti
    \item Validazione dei tipi
    \item Help auto-generato
    \item Composizione di comandi (gruppi)
\end{itemize}

\section{Comandi Disponibili}

\subsection{parse}

Elabora file XML e li converte in Parquet:

\begin{lstlisting}[language=bash]
docker compose run --rm etl python -m src.cli parse \
    --input data/ \
    --output public/parquet \
    --workers 8
\end{lstlisting}

\begin{table}[h]
\centering
\begin{tabular}{@{}llp{5cm}@{}}
\toprule
\textbf{Opzione} & \textbf{Default} & \textbf{Descrizione} \\
\midrule
\texttt{-i, --input} & (required) & File o directory input \\
\texttt{-o, --output} & \texttt{public/parquet} & Directory output \\
\texttt{-w, --workers} & 4 & Numero processi paralleli \\
\bottomrule
\end{tabular}
\caption{Opzioni comando parse}
\end{table}

\subsection{query}

Esegue query SQL interattive via DuckDB:

\begin{lstlisting}[language=bash]
# Query semplice
docker compose run --rm etl python -m src.cli query \
    --table aiuti --limit 5

# Query SQL custom
docker compose run --rm etl python -m src.cli query \
    --table strumenti \
    --query "SELECT ANNO, SUM(IMPORTO_NOMINALE) as tot 
             FROM read_parquet('public/parquet/strumenti/**/*.parquet') 
             GROUP BY ANNO ORDER BY ANNO"
\end{lstlisting}

\begin{table}[h]
\centering
\begin{tabular}{@{}llp{5cm}@{}}
\toprule
\textbf{Opzione} & \textbf{Default} & \textbf{Descrizione} \\
\midrule
\texttt{-t, --table} & (required) & Tabella target \\
\texttt{-q, --query} & — & Query SQL custom \\
\texttt{-l, --limit} & 10 & Limite risultati \\
\bottomrule
\end{tabular}
\caption{Opzioni comando query}
\end{table}

\subsection{export}

Esporta tabelle in formato CSV o TXT:

\begin{lstlisting}[language=bash]
docker compose run --rm etl python -m src.cli export \
    --table aiuti \
    --format csv \
    --output public/exports/aiuti.csv \
    --delimiter ","
\end{lstlisting}

\subsection{export-aggregated}

Export specializzato con join e aggregazioni pre-calcolate:

\begin{lstlisting}[language=bash]
docker compose run --rm etl python -m src.cli export-aggregated \
    --output public/exports/aggregated.csv
\end{lstlisting}

Questo comando:
\begin{enumerate}
    \item Effettua join tra AIUTI, COMPONENTI e STRUMENTI
    \item Aggrega \texttt{IMPORTO\_NOMINALE} e \texttt{ELEMENTO\_DI\_AIUTO}
    \item Conta componenti e strumenti per aiuto
    \item Produce un file CSV per anno per ottimizzare la memoria
\end{enumerate}

\section{Progress Bar}

Durante l'elaborazione, \texttt{tqdm} fornisce feedback visivo in tempo reale:

\begin{lstlisting}[language=bash, caption=Output con progress bar]
Found 165 XML files to process
Starting processing with 8 workers...
Processing files: 100%|===================| 165/165 [02:15<00:00]
Processing completed in 135.42 seconds
Total processed records: {'aiuti': 2847293, 'componenti': ...}
\end{lstlisting}

\section{Error Handling}

Il sistema gestisce gli errori in modo graceful:

\begin{itemize}
    \item \textbf{File corrotti}: loggati ma non bloccano l'esecuzione
    \item \textbf{failures.txt}: lista dei file falliti salvata per retry
    \item \textbf{Statistiche accurate}: solo file processati correttamente conteggiati
\end{itemize}
